\documentclass[11pt]{article}
\usepackage[a4paper,margin=28mm]{geometry}
\usepackage{amsmath,amssymb,amsthm}
\usepackage[T1]{fontenc}
\usepackage{lmodern}
\usepackage{microtype}
\usepackage{hyperref}
\hypersetup{colorlinks=true,linkcolor=blue,citecolor=blue,urlcolor=blue,
  pdftitle={Failure of exterior observability at equal bulk and surface wave speeds in an annulus},
  pdfauthor={Kaiwen Jin}}
\newtheorem{theorem}{Theorem}
\newtheorem{lemma}{Lemma}
\newcommand{\Ai}{\operatorname{Ai}}
\newcommand{\Bi}{\operatorname{Bi}}
\title{Failure of exterior observability at equal bulk and surface wave speeds in an annulus}
\author{Kaiwen Jin\\\small\href{mailto:yui.ui.siki@gmail.com}{yui.ui.siki@gmail.com}}
\date{}
\begin{document}
\maketitle

\begin{abstract}
We consider a wave equation in a planar annulus, with a Dirichlet condition on
the outer circle and a dynamic wave equation on the inner circle. The bulk and
surface propagation speeds are both one. For every finite observation time,
we construct smooth real solutions with unit initial energy whose exterior
normal-derivative observation tends to zero. The construction uses exact
separated solutions with angular order \(m\) and frequency
\(\lambda_m=m+O(m^{-1/3})\). Large-order Bessel asymptotics at the inner turning
point give a lower bound for the surface energy, while an exact Wronskian
identity fixes the outer flux. Thus exterior observability fails at equality
of the two speeds in this annulus.
\end{abstract}

\section{The model and the result}

Dynamic boundary conditions describe a vibrating surface coupled to a wave
in the surrounding medium. We ask whether measurements on the exterior
boundary determine the entire bulk and surface energy. Let
\[
\Omega=\{x\in\mathbb R^2:1<|x|<2\},\qquad
\Gamma_1=\{|x|=1\},\qquad \Gamma_0=\{|x|=2\}.
\]
The equal-speed system is
\begin{equation}\label{eq:model}
\begin{aligned}
z_{tt}-\Delta z&=0&&\text{in }\Omega,\\
z&=0&&\text{on }\Gamma_0,\\
z_\Gamma&=z|_{\Gamma_1},\\
z_{\Gamma,tt}-\Delta_\Gamma z_\Gamma+\partial_\nu z&=0
&&\text{on }\Gamma_1 .
\end{aligned}
\end{equation}
Here \(\nu\) is the outward normal of the annulus and \(\Delta_\Gamma\)
is the Laplace operator of the unit circle. In polar coordinates,
\(\partial_\nu=-\partial_r\) on \(\Gamma_1\),
\(\partial_\nu=\partial_r\) on \(\Gamma_0\), and
\(\Delta_\Gamma=\partial_{\theta\theta}\). The energy and observation are
\begin{align}
E_z(t)&=\int_\Omega\bigl(|z_t|^2+|\nabla z|^2\bigr)\,dx
+\int_{\Gamma_1}\bigl(|z_{\Gamma,t}|^2+
|\nabla_\Gamma z_\Gamma|^2\bigr)\,dS,\label{eq:energy}\\
\mathcal O_T(z)&=\int_0^T\int_{\Gamma_0}|\partial_\nu z|^2\,dS\,dt.
\label{eq:observation}
\end{align}
All solutions constructed below are real and smooth up to both boundaries
and for all real times.

Baudouin, Dard\'e, Ervedoza and Mercado study the annular system with
surface mass \(\alpha\) and tangential coefficient \(\beta\)
\cite[Eq.~(1.3), pp.~60--61]{BDEM}. Their Theorem~2.2 treats
\(\alpha<\beta\), while Theorem~2.4 gives failure of observability for
\(\alpha>\beta\). Neither theorem includes \(\alpha=\beta=1\), which is
\eqref{eq:model}. The review of Chorfi, Maniar and Morales
\cite[Sec.~5.1]{CMM} explicitly leaves the equal-speed case open.
The result here answers the concrete annular question in
\cite{AIM385} negatively.

\begin{theorem}\label{thm:main}
There is an infinite set \(\mathcal M\) of positive integers and smooth real
solutions \(w_m\) of \eqref{eq:model}, \(m\in\mathcal M\), such that
\[
E_{w_m}(0)=1,\qquad
\mathcal O_T(w_m)\le C T m^{-1/3}
\quad\text{for every finite }T>0,
\]
where \(C\) is independent of \(m\) and \(T\).
Consequently, for no finite \(T>0\) does a finite constant \(C_T\) satisfy
\(E_z(0)\le C_T\mathcal O_T(z)\) for all smooth compatible solutions.
In particular, no finite threshold time \(T_0\) has this property for every
\(T>T_0\).
\end{theorem}

We construct exact modes rather than approximate solutions. Their radial
equation has a turning point near the inner circle when \(\lambda\) is close
to \(m\). The inner boundary equation can be solved with a frequency
correction of order \(m^{-1/3}\). On a subsequence, the resulting surface
kinetic energy grows while the unnormalized exterior flux remains fixed.

\section{An exact family of radial modes}

Let \(m\) be a positive integer and \(J_m,Y_m\) the conventional Bessel functions of order \(m\).
For \(\lambda>0\), the radial equation for angular dependence
\(\cos(m\theta)\) is
\begin{equation}\label{eq:radial}
f''+r^{-1}f'+(\lambda^2-m^2/r^2)f=0 .
\end{equation}
To impose the outer Dirichlet condition, define
\begin{equation}\label{eq:F}
F_m(r;\lambda)=J_m(\lambda r)Y_m(2\lambda)
-Y_m(\lambda r)J_m(2\lambda),\qquad 1\le r\le2 .
\end{equation}
The Bessel equation gives \eqref{eq:radial} and \(F_m(2;\lambda)=0\).
The Wronskian identity
\cite[Eq.~10.5.2]{DLMF}
\[
J_m(x)Y_m'(x)-J_m'(x)Y_m(x)=\frac{2}{\pi x}
\]
also gives the exact radial derivative
\begin{equation}\label{eq:flux}
\partial_rF_m(2;\lambda)
=\lambda\bigl[J_m'(2\lambda)Y_m(2\lambda)
-Y_m'(2\lambda)J_m(2\lambda)\bigr]=-\frac1\pi .
\end{equation}
In particular \(F_m\) is nonzero. Bessel primes mean argument derivatives;
radial derivatives are written as \(\partial_rF_m\).

The inner dynamic condition for a mode
\(\sin(\lambda t)F_m(r;\lambda)\cos(m\theta)\) is exactly
\begin{equation}\label{eq:G}
G_m(\lambda):=\partial_rF_m(1;\lambda)
+(\lambda^2-m^2)F_m(1;\lambda)=0.
\end{equation}
Indeed the inner outward normal is \(-\partial_r\), so its boundary
equation is \((m^2-\lambda^2)F_m(1;\lambda)
-\partial_rF_m(1;\lambda)=0\). It remains to find roots near \(\lambda=m\)
with a sufficiently large inner value.

\begin{lemma}\label{lem:shift}
Set
\[
\kappa=\sqrt3-\frac{\pi}{3},\quad
\phi_m=m\kappa-\frac{\pi}{4},\quad
A=\left(\frac{2}{\pi\sqrt3}\right)^{1/2},
\]
Let \(\Ai,\Bi\) be the conventional Airy functions, and put
\[
\begin{aligned}
a&=2^{1/3}\Ai(0)>0,&\qquad b&=2^{1/3}\Bi(0)>0,\\
p&=-2^{2/3}\Ai'(0),& q&=2^{2/3}\Bi'(0).
\end{aligned}
\]
For every fixed \(M>0\), as positive integers \(m\) tend to infinity,
the following estimates hold uniformly for
\(\lambda=m+s m^{-1/3}\), \(|s|\le M\):
\begin{align}
J_m(\lambda)&=a m^{-1/3}+o(m^{-1/3}),&
Y_m(\lambda)&=-b m^{-1/3}+o(m^{-1/3}),\label{eq:inner-values}\\
J_m'(\lambda)&=p m^{-2/3}+o(m^{-2/3}),&
Y_m'(\lambda)&=q m^{-2/3}+o(m^{-2/3}),\label{eq:inner-derivatives}\\
J_m(2\lambda)&=A m^{-1/2}(\cos\phi_m+o(1)),&
Y_m(2\lambda)&=A m^{-1/2}(\sin\phi_m+o(1)).\label{eq:outer-values}
\end{align}
\end{lemma}

\begin{proof}
We first use standard expansions only at the fixed arguments \(m\) and
\(2m\), and then prove the required uniformity from the Bessel equation.
In the uniform Airy expansions
\cite[Eqs.~10.20.4, 10.20.5, 10.20.7, 10.20.8]{DLMF}, the Bessel
argument is \(x=m\eta\), with normalized argument \(\eta>0\).
The value formulas and their separately stated argument-derivative
formulas, evaluated at \(\eta=1\), give \eqref{eq:inner-values} and
\eqref{eq:inner-derivatives} at \(\lambda=m\).
To check the prefactors, the defining turning variable satisfies
\(\zeta(\eta)=2^{1/3}(1-\eta)+O((1-\eta)^2)\); hence the continuous limits of
\((4\zeta/(1-\eta^2))^{1/4}\) and
\((2/\eta)((1-\eta^2)/(4\zeta))^{1/4}\) are \(2^{1/3}\) and
\(2^{2/3}\). The leading coefficient functions in the value and
argument-derivative expansions equal one, and the remaining coefficients
are smooth at \(\zeta=0\)
\cite[Eqs.~10.20.10--10.20.14]{DLMF}.
The Airy values and signs are given in
\cite[Eqs.~9.2.3--9.2.6]{DLMF}.

The fixed-angle Debye expansions
\cite[Eqs.~10.19.5--10.19.7]{DLMF} use the argument \(x=m\sec\beta\)
and phase \(m(\tan\beta-\beta)-\pi/4\), where the angle
\(\beta\in(0,\pi/2)\) is fixed. At \(\beta=\pi/3\), the argument is
\(2m\) and the phase is \(\phi_m\), so they give
\[
\begin{aligned}
J_m(2m)&=A m^{-1/2}\cos\phi_m+O(m^{-3/2}),\\
Y_m(2m)&=A m^{-1/2}\sin\phi_m+O(m^{-3/2}),\\
|J_m'(2m)|+|Y_m'(2m)|&=O(m^{-1/2}).
\end{aligned}
\]
These error estimates are absolute, including at zeros of the leading
trigonometric factors.

For either \(y=J_m\) or \(y=Y_m\), its equation is
\[
y''(x)=-x^{-1}y'(x)-(1-m^2/x^2)y(x).
\]
Put \(h=M m^{-1/3}\). On \(|x-m|\le h\), for all sufficiently large \(m\),
\[
x^{-1}\le2/m,\qquad |1-m^2/x^2|\le C_M m^{-4/3}.
\]
Let \(U=\sup|y|\) and \(V=\sup|y'|\) on this interval. Integrating from
\(m\), on either side, gives
\[
U\le |y(m)|+hV,\qquad
V\le |y'(m)|+h(2V/m+C_M m^{-4/3}U).
\]
Substitute the first bound into the second. The coefficient of \(V\) on
the right tends to zero and can be absorbed. The endpoint scales give
\(V=O_M(m^{-2/3})\) and \(U=O_M(m^{-1/3})\). Integrating once more yields
\[
y(x)-y(m)=O_M(m^{-1}),\qquad
y'(x)-y'(m)=O_M(m^{-2}).
\]
These changes are smaller than the respective leading scales.

On \(|x-2m|\le2h\), both equation coefficients are bounded. If
\(W=\sup(|y|+|y'|)\) there, integration gives
\[
W\le |y(2m)|+|y'(2m)|+C(2h)W.
\]
Since \(h\to0\), absorption gives \(W=O_M(m^{-1/2})\), and
\(y(x)-y(2m)=O_M(m^{-5/6})=o(m^{-1/2})\). This proves all the
uniform shifted estimates without differentiating an asymptotic remainder.
\end{proof}

\begin{lemma}\label{lem:roots}
There exist constants \(M,\delta>0\), an infinite set
\(\mathcal M\) of positive integers, and exact roots
\(\lambda_m>0\) of \eqref{eq:G} such that, for all sufficiently large
\(m\in\mathcal M\),
\[
|\lambda_m-m|\le M m^{-1/3},\qquad
|F_m(1;\lambda_m)|\ge \frac{A\delta}{2}m^{-5/6}.
\]
\end{lemma}

\begin{proof}
Define the two phase factors
\[
S_m=a\sin\phi_m+b\cos\phi_m,\qquad
B_m=p\sin\phi_m-q\cos\phi_m .
\]
At the inner circle, the determinant \eqref{eq:F} and its radial
derivative read
\[
\begin{aligned}
F_m(1;\lambda)&=J_m(\lambda)Y_m(2\lambda)-Y_m(\lambda)J_m(2\lambda),\\
\partial_rF_m(1;\lambda)
&=\lambda\bigl[J_m'(\lambda)Y_m(2\lambda)-Y_m'(\lambda)J_m(2\lambda)\bigr].
\end{aligned}
\]
Lemma~\ref{lem:shift} therefore gives, uniformly on each fixed \(|s|\le M\),
\begin{align}
F_m(1;m+s m^{-1/3})&=A m^{-5/6}(S_m+o(1)),\label{eq:Finner}\\
\partial_rF_m(1;m+s m^{-1/3})
&=A m^{-1/6}(B_m+o(1)).\label{eq:Frinner}
\end{align}
The second formula includes the factor \(\lambda/m=1+o(1)\) from
radial differentiation.

We choose a subsequence where the inner value does not cancel. Write
\(S_m=R\sin(\phi_m+\vartheta)\), where
\(R=(a^2+b^2)^{1/2}>0\). Since \(0<\kappa<\pi\),
\[
\sin\kappa=\sin(x+\kappa)\cos x-\cos(x+\kappa)\sin x
\]
implies \(|\sin x|+|\sin(x+\kappa)|\ge\sin\kappa>0\).
Thus in each disjoint pair \(m=2k,2k+1\), \(k\ge1\), at least one
index satisfies
\[
|S_m|\ge\delta:=\frac{R\sin\kappa}{2}.
\]
Select one such index from each pair to obtain the infinite set
\(\mathcal M\). This argument needs no equidistribution of the phases.

The exact expansion
\(\lambda^2-m^2=2s m^{2/3}+s^2m^{-2/3}\), together with
\eqref{eq:Finner}--\eqref{eq:Frinner}, gives
\begin{equation}\label{eq:Gscale}
m^{1/6}G_m(m+s m^{-1/3})
=A\bigl(B_m+2sS_m+o(1)\bigr)
\end{equation}
uniformly for each fixed \(|s|\le M\). Choose \(M\) once and for all with
\[
2M\delta>|p|+|q|+2.
\]
For sufficiently large \(m\in\mathcal M\), the right side of
\eqref{eq:Gscale} has opposite signs at \(s=-M\) and \(s=M\):
the terms \(\mp2MS_m\) dominate \(B_m\) and the uniform remainder.
Both frequency endpoints are positive. The intermediate value theorem
therefore supplies an exact root \(\lambda_m\) in that interval.
Finally, \eqref{eq:Finner} and \(|S_m|\ge\delta\) give the asserted
inner-value bound throughout the same interval, in particular at its root.
\end{proof}

\section{Energy and exterior observation}

\begin{proof}[Proof of Theorem~\ref{thm:main}]
Discard the finitely many low orders excluded in Lemma~\ref{lem:roots},
retaining the notation \(\mathcal M\). Put
\[
f_m(r)=F_m(r;\lambda_m),\qquad
z_m(t,r,\theta)=\sin(\lambda_m t)f_m(r)\cos(m\theta).
\]
The polar Laplacian and \eqref{eq:radial} give the bulk wave equation.
The outer trace is zero, and \eqref{eq:G} gives the inner equation
with the stated outward normal. Integer angular order ensures periodicity.
Every Bessel argument is positive, so the modes are real and smooth on the
closed annulus for all real times. These exact identities also give every
initial compatibility condition.

At \(t=0\), displacement is zero and velocity is
\(\lambda_m f_m(r)\cos(m\theta)\). Since
\(\int_0^{2\pi}\cos^2(m\theta)\,d\theta=\pi\), the full energy
\eqref{eq:energy} is
\begin{equation}\label{eq:modeenergy}
E_{z_m}(0)=\pi\lambda_m^2
\left(\int_1^2|f_m(r)|^2r\,dr+|f_m(1)|^2\right).
\end{equation}
For all sufficiently large selected \(m\), \(\lambda_m\ge m/2\).
The surface term alone and Lemma~\ref{lem:roots} yield
\begin{equation}\label{eq:energylower}
E_{z_m}(0)\ge\frac{\pi A^2\delta^2}{16}m^{1/3}>0 .
\end{equation}
On the outer circle \(dS=2\,d\theta\), and \eqref{eq:flux} gives
\begin{equation}\label{eq:modeobservation}
\mathcal O_T(z_m)=\frac{2}{\pi}\int_0^T\sin^2(\lambda_m t)\,dt
\le\frac{2T}{\pi}.
\end{equation}
Normalize by \(w_m=z_m/\sqrt{E_{z_m}(0)}\). Energy and observation are
quadratic, so \eqref{eq:energylower}--\eqref{eq:modeobservation} give
\[
E_{w_m}(0)=1,\qquad
\mathcal O_T(w_m)\le\frac{32T}{\pi^2 A^2\delta^2}m^{-1/3}
\longrightarrow0
\]
for every fixed finite \(T>0\). A finite \(C_T\) in the proposed inequality
would require \(1\le C_T\mathcal O_T(w_m)\) for this sequence, which is
impossible. Choosing any finite positive \(T>T_0\) also excludes every
finite threshold \(T_0\).
\end{proof}

The conclusion concerns precisely the unit-speed system on the annulus
with radii one and two. The estimate above establishes failure of energy
observability; it does not identify an optimal weaker observation norm or
a stabilization rate.

\paragraph{AI-use statement.}
OpenAI Codex assisted with literature retrieval, the spectral construction,
checking the argument, and drafting this manuscript. The conclusion is
supported by the analytic proof above; no numerical calculation is used as
a proof premise. Independent mathematical review is requested.

\begin{thebibliography}{9}
\bibitem{BDEM}
L. Baudouin, J. Dard\'e, S. Ervedoza and A. Mercado,
\emph{A unified strategy for observability of waves in an annulus with
various boundary conditions},
Mathematical Reports \textbf{24(74)} (2022), 59--112.
\url{https://imar.ro/journals/Mathematical_Reports/Pdfs/2022/1-2/5.pdf}.

\bibitem{CMM}
S. E. Chorfi, L. Maniar and R. Morales,
\emph{Controllability and Inverse Problems for Hyperbolic and Dispersive
Equations with Dynamic Boundary Conditions},
arXiv:2505.14795v1, 20 May 2025.
\url{https://arxiv.org/abs/2505.14795v1}.

\bibitem{DLMF}
NIST Digital Library of Mathematical Functions,
Sections 9.2, 10.2, 10.5, 10.19, 10.20 and 10.41.
\url{https://dlmf.nist.gov/}.
Consulted 12 October 2026.

\bibitem{AIM385}
MColbrook/AIM, \emph{Observability of an acoustic shell
at equal bulk and surface speeds}, AIM problem 385, statement at revision
\texttt{c9929805cece044705abaaebe1cebddac1ba11a6}.
\href{https://github.com/MColbrook/AIM/blob/c9929805cece044705abaaebe1cebddac1ba11a6/problems/385-equal-speed-wentzell-wave-observability.md}{Problem statement at the cited revision}.
\end{thebibliography}
\end{document}
